\section{Empirical Testability via Revealed Sacrifice}
\label{sec:empirical}

The quantification thesis of Section~\ref{sec:lineage} is
empirically operational because the
\citep[Revealed~Sacrifice]{lasser2026paper8a} channel maps
directly onto observable economic ledgers.  Purchase records,
time-use surveys, labour-supply panels, and public-commitment
ledgers are revealed-sacrifice data in the GFM sense; the
\citep[Need~Sufficiency]{lasser2026paper8b} diagnostic
architecture (gap decomposition, wireheading-consistent HHI,
need-access cost) is the interpretation layer that turns raw
observation into testable predictions.  This section catalogues
the formal instruments and operationalizes the Goodhart theorem
of Section~\ref{sec:goodhart} as a cross-domain empirical test.
The worked examples of Section~\ref{sec:worked_examples} then
illustrate the instruments concretely.

A note on the relationship to existing operational welfare
indices.  The Human Development
Index~\citep{anand1994human}, the OECD Better Life Index, the
UNDP Multidimensional Poverty Index, and similar
deployment-readable welfare instruments are operational
measurement tools rather than formal frameworks: they aggregate
chosen dimensions through chosen weighting schemes to produce
deployment-readable scalars.  The instruments below are
\emph{formal analogues} to these operational indices and supply
an interpretation layer for analogous index choices: deployments
using the operational indices can read the indices' implicit
individuation and weighting decisions through the GFM apparatus
to identify which channel of Section~\ref{sec:lineage} is
operative when the index moves.  The match between a specific
operational index and a specific GFM instrument is not exact
in general (the operational indices were not designed to be
GFM-compatible), but the structural correspondence is enough
to position the operational indices and the formal instruments
as complementary rather than competing.

\subsection{Instruments: GFM machinery as economic-data interpretation}
\label{sec:instruments}

Six formal instruments interpret revealed-sacrifice data.  Each
maps a familiar economic data source to a specific component of
the formalization, and each is what one or more channels of
Section~\ref{sec:lineage} act on.

\begin{table}[ht]
\centering
\footnotesize
\setlength{\tabcolsep}{4pt}
\begin{tabular}{>{\raggedright\arraybackslash}p{0.20\linewidth}
                >{\raggedright\arraybackslash}p{0.27\linewidth}
                >{\raggedright\arraybackslash}p{0.27\linewidth}
                >{\raggedright\arraybackslash}p{0.18\linewidth}}
\toprule
\textbf{Instrument} & \textbf{Economic data source} &
\textbf{What it measures} & \textbf{Channel(s)} \\
\midrule
Two channels (money + time) &
Purchase ledgers + time-use surveys &
Independent observation modalities of revealed sacrifice &
Channel~1 \\
\addlinespace
Aggregate lower bound $\betaL$ &
Trade-event ledger restricted to admissible events &
Diagnostic ceiling $(1-\betaL)/\betaL$ on $\epsnonres$ &
Channels~1, 2 \\
\addlinespace
Five-cell $\delta$-decomposition &
Cell-wise attribution across covered, dormant, restricted, residual, boundary &
Where the gap lives across structural categories &
Selector for Channels~1--4 \\
\addlinespace
Wireheading-consistent HHI &
Population-scale trade-flow concentration &
Concentration consistent with Conjecture~1 &
Diagnostic for Channel~1 saturation \\
\addlinespace
Need-access cost (NAC) &
Below-sufficiency expenditure ledger &
Need-side share of revealed sacrifice &
Channels~3, 4 \\
\addlinespace
Per-event $\alpha_{n,c}$ &
Bundle decomposition extension &
Per-capability attribution within bundles &
Channel~4 \\
\bottomrule
\end{tabular}
\caption{The six formal instruments connecting GFM machinery to
economic data sources.  Each instrument operationalizes one or
more components of $\epsgap$ and is what at least one channel of
Section~\ref{sec:lineage} acts on.  The five-cell
$\delta$-decomposition plays a selector role across all four
channels, mirroring the diagnostic role the SSF~2009
decomposition plays in welfare-economics measurement
methodology.}
\label{tab:instruments}
\end{table}

\paragraph{Two channels (money + time).}
\citep[Section~4]{lasser2026paper8a} establishes two
substantively distinct observation modalities: money sacrifice
(priced trade events) and time sacrifice (foregone-wage
opportunity costs).  The two channels triangulate the same
underlying realized exercise; cross-channel agreement is a
robustness check, and divergences flag attestation or
individuation issues.  Predictions~\ref{pred:bundle_completion}
through~\ref{pred:saturation} of
Section~\ref{sec:structural_predictions} are testable through
either channel, with cross-channel triangulation providing
independent confirmation.

\paragraph{Aggregate lower bound $\betaL$.}
\citep[Theorem~2]{lasser2026paper8a} bounds the realized
exercise volume from below by an aggregate of admissible
revealed-sacrifice events.  $\betaL$ is the aggregate lower
ratio; $(1 - \betaL)/\betaL$ is the corresponding
deployment-readable diagnostic ceiling on the within-channel
disagreement gap $\epsnonres$
(Remark~\ref{rem:agg_vs_sup}).  Computing $\betaL$ requires:
(i)~a defined trade window~$W$, (ii)~classification of events
as admissible or not under~S0, S1, S4(a), and the
genuine-exchange conditions, and (iii)~aggregation under
Part~A (bundle-level) or Part~B (per-capability) modes.

\paragraph{Five-cell $\delta$-decomposition.}
\citep[Definition~8]{lasser2026paper8b} attributes the gap
across $\delta_{\mathrm{covered}}$ (channel observed),
$\delta_{\mathrm{dormant}}$ (channel could reach but hasn't),
$\delta_{\mathrm{restricted}}$ (deliberately non-exercised),
$\delta_{\mathrm{residual}}$ ($\ResS$ proper), and
$\delta_{\mathrm{boundary}}$ (boundary-residual).  The
attribution is a structure question (``where does the gap
live?'') and is the deployment instrument that selects which
channel of Section~\ref{sec:lineage} binds for the given gap
profile.

\paragraph{Wireheading-consistent HHI.}
\citep[Propositions~6~and~7]{lasser2026paper8b} construct a
Schur-convex concentration measure on trade-flow distributions
across bundle categories.  The HHI is the population-scale
empirical witness for Conjecture~\ref{conj:optimization_pressure}: high HHI
indicates concentration consistent with optimization-pressure
exploitation of $P$-$T$ divergence regions, though as
Remark~\ref{rem:hhi_empirical_witness} insists, the witness is
correlational rather than causal.

\paragraph{Need-access cost (NAC).}
\citep[Definition~9]{lasser2026paper8b} tracks below-sufficiency
expenditure: revealed sacrifice on capabilities below the
polarity boundary, where retention is degrading and the agent
has no non-degrading alternative.  NAC distinguishes need-share
from want-share contributions to revealed sacrifice and is the
operational substrate for Prediction~\ref{pred:mixed_polarity}.

\paragraph{Per-event $\alpha_{n,c}$.}
\citep[Definition~4]{lasser2026paper8a} defines the
event-local bundle decomposition coefficients $\alpha_{n,c}$
that partition a bundle's contribution to $\Delta\volL(X_n)$
across capability components.  Per-event decomposition is the
substrate for Channel~4 and for
Predictions~\ref{pred:bundle_completion}
and~\ref{pred:coop_premium}.

\subsection{Operationalizing the Goodhart bound}
\label{sec:operationalizing}

The Goodhart theorem of Section~\ref{sec:goodhart} is testable
cross-domain and longitudinally.  Both tests follow the same
structural recipe but ask different questions of the data.

\paragraph{Cross-domain test.}
For a collection of capability domains $\mathcal{D}_1, \dots,
\mathcal{D}_m$ in which the framework operates, compute the
diagnostic ceiling $(1 - \betaL)/\betaL$ on each domain's
$\epsnonres$ and the cell-wise attribution via the
$\delta$-decomposition.  For a fixed alignment property $g$
(say, anti-monopolar threshold accuracy), measure $g$'s
deviation from a domain-specific operational truth.
Theorem~\ref{thm:goodhart} certifies that domains with
smaller diagnostic-ceiling values support smaller \emph{certified
upper bounds} on $|g(T) - g(P)|$; the cross-domain empirical
test is therefore a directional-consistency check between the
diagnostic ceiling and observed deviations, with the cell-wise
attribution telling the deployer which channel intervention
would produce the within-domain tightening of the certified
bound.  This test does not require closing the per-property
Lipschitz constants; it requires only checking the directional
consistency between cross-domain $\epsnonres$ and cross-domain
$|g(T) - g(P)|$.  A domain
violating the directional consistency is either a Lipschitz
failure (the property's sensitivity to $\epsnonres$ is
non-Lipschitz on the domain) or an attestation failure (the
$\betaL$ ceiling is unreliable for the domain).

\paragraph{Longitudinal test.}
For a single deployment over a window $W = [t_0, t_0 + T]$ in
which the deployment performs an admissible intervention on one
of the four channels (more events, sharper attestation,
re-individuation, bundle decomposition), measure the
intervention's effect on $\betaL$, on the cell-wise attribution,
and on $g$'s deviation from the operational truth.
Theorem~\ref{thm:goodhart} certifies a monotone tightening of
the \emph{upper bound} on $|g(T) - g(P)|$ with the
intervention's effect on $\epsnonres$; the empirical test is the
directional-consistency check on observed deviations.  Failure
of monotonicity in the certified bound is either a Lipschitz
violation, an attestation regression, or an unaccounted-for
counter-channel effect (e.g., re-individuation that increases
within-channel disagreement at the same time as it decreases
the floor; cf.\ Remark~\ref{rem:floor_binds}).

\paragraph{Distinguishing test outcomes.}
Three distinct test outcomes have distinct interpretations:
\begin{enumerate}
\item \emph{Tightening as predicted, attribution consistent
  with intervention.}  The framework's prediction is supported,
  and the deployment can plan further interventions on the
  binding channel.
\item \emph{Tightening as predicted, attribution inconsistent
  with intervention.}  The intervention worked, but for
  different reasons than expected.  This is informative: it
  suggests the cell-wise attribution is mis-specified (the
  intervention reduced the wrong cell's gap), which is a
  diagnostic-decomposition issue rather than a Goodhart
  issue.
\item \emph{No tightening or anti-tightening.}  The intervention
  failed.  The most common failure modes are: (a)~Lipschitz
  violation (the alignment property is not Lipschitz on the
  relevant capability subspace under the deployment's
  individuation), (b)~adversarial attestation degradation
  (Open~Question~8 of Section~\ref{sec:discussion}), and
  (c)~cross-channel offset (one channel improved while another
  worsened).  Distinguishing the failure modes requires
  cell-wise attribution before and after the intervention.
\end{enumerate}

\paragraph{Sup-norm coarseness and empirical sensitivity.}
A caveat for the empirical programme that does not affect the
theory.  The sup-norm definition of $\epsnonres$
(Definition~\ref{def:eps_nonres}) is conservative
(Remark~\ref{rem:supnorm_conservative}): a deployment that
tightens $\epsnonres$ on 99\,\% of capability dimensions while
slightly worsening it on one shows no improvement under the
sup-norm.  How this caveat propagates to the diagnostic
ceiling $(1-\betaL)/\betaL$ depends on how the deployment
reads the ceiling.  When the deployment allocates per-capability
ceiling components under
Assumption~\ref{ass:per_cap_admissibility} part~(i) and then
takes the worst-case across them as a sup-norm proxy, the
diagnostic ceiling inherits the worst-dimension conservatism;
when the deployment reads the ceiling purely as an aggregate
(without per-capability allocation), the aggregate may still
register average progress even if a single dimension regresses,
but the bound on the sup-norm $\epsnonres$ is then no longer
tight (Remark~\ref{rem:agg_vs_sup}).  In practice the
cross-domain and longitudinal tests above may therefore be
\emph{insensitive to genuine progress on most dimensions when a
single dimension is regressing} (under the sup-norm reading
that bounds $\epsnonres$ tightly), and a deployment may be
meaningfully improving without the sup-norm-bounded diagnostic
registering it.  This is a feature for safety-conservative
alignment monitoring (one badly-proxied dimension is
sufficient to produce a misaligned decision), but it is a
limitation for deployment-progress measurement, where weighted
norms or per-dimension diagnostics may be operationally more
informative.  Open~Question~\ref{oq:nonmonotonicity} flags the
weighted-norm direction as a candidate refinement; the
empirical operator should interpret a flat sup-norm-bounded
diagnostic under intervention as ``no progress on the worst
dimension'' rather than as ``no progress at all,'' and may
choose to track the aggregate $\betaL$ separately as a
deployment-progress signal that the sup-norm bound's
conservatism does not eat.

\subsection{Testing the four channels}
\label{sec:testing_channels}

Each of the four channels of Section~\ref{sec:lineage} has a
specific test that the instruments support.

\paragraph{Channel 1 (observation density).}
The longitudinal trajectory of $\betaL$ under admissible
accumulation should be non-decreasing
(\citep[Propositions~2~and~3]{lasser2026paper8a}).  Concretely:
add a sequence of admissible trade events to a deployment's
ledger one at a time; $\betaL$ at each step should be at least
as large as the previous step.  A decrease falsifies either the
admissibility classification or the framework's
non-decreasing claim.  In practice, deployment data is added
in batches; batch-wise non-decreasing is the operational
version of the test.

\paragraph{Channel 2 (attestation quality).}
The attested-share of admissible events (events that pass the
S0, S1, S4(a) checks under the deployment's institutional
attestation infrastructure) should rise as the attestation
layer matures.  The diagnostic-ceiling
$(1 - \betaL)/\betaL$ at fixed observation count should
tighten as the attested share rises (events previously
classified as non-admissible become admissible under sharper
attestation).  Falsification: $\betaL$ does not respond to
attested-share improvements, indicating that the attestation
layer's classification is not improving the operational
admissibility set or that the events being attested are
non-binding.

\paragraph{Channel 3 (individuation discipline).}
A re-individuation move (e.g., treating household labour as a
distinct capability with imputed-time pricing) should reduce
$\epsfloor$ and add a corresponding contribution to
$\epsnonres$ from the newly-admitted capability.  The net
effect on $\epsgap = \epsnonres + \lambda \cdot \epsfloor$
depends on whether the new capability's attested share is
high (in which case its $\epsnonres$ contribution is small,
and the net move is a tightening) or low (in which case the
move is neutral or even alignment-loosening, with the
$\lambda$-weighted net $\Delta_3^{\lambda}$ approaching zero
or going negative; cf.\ Section~\ref{sec:channel3_limitation}).
Test: pre/post comparison of $\epsfloor$, $\epsnonres$, and
the alignment property; the move tightens iff the new
capability's attested share exceeds a deployment-specific
threshold determined by $\lambda$.

\paragraph{Channel 4 (bundle decomposition).}
The per-event $\alpha_{n,c}$ decomposition (when the
deployment's commitment scheme supports it) should produce
within-bundle component-wise lower bounds that, when
aggregated, are at least as tight as the bundle-level Part-A
aggregation of \citep[Theorem~2]{lasser2026paper8a}.
Operationally: compute $\betaL$ under both Part-A and Part-B
modes for the same trade-event set; Part-B (per-capability)
should be at least as informative as Part-A on the
within-bundle poset-independence subset of events.  If
Part-B is tighter, Channel~4 is delivering its predicted
tightening; if not, the bundle-decomposition extension is
either inactive or applied to events outside its
admissibility scope.

\subsection{Methodology template and forward}
\label{sec:methodology_template}

\citep[Section~6]{lasser2026paper8b} works through a four-phase
end-to-end demonstration of the
\citep[Need~Sufficiency]{lasser2026paper8b} diagnostic
architecture on a single domain: capability inventory $\to$
trade-event ledger $\to$ cell-wise attribution $\to$
deployment intervention.  That worked example is a methodology
template.  The four worked examples of
Section~\ref{sec:worked_examples} below illustrate the full
machinery of the present paper at four scales: (i)~micro-level
divergence between a GFM prediction and the utility-theoretic
account (Section~\ref{sec:wex_boat}), (ii)~the Goodhart
gradient under variable observation infrastructure
(Section~\ref{sec:wex_goodhart_gradient}),
(iii)~fungibility-collapse limit at macro scale
(Section~\ref{sec:wex_collapse}), and (iv)~mixed-polarity
bundle-for-bundle trades with deployment misalignment
consequences (Section~\ref{sec:wex_relocation}).  Each example
exercises specific instruments from Table~\ref{tab:instruments}
and demonstrates concrete tightening predictions or failure
modes; the four examples together cover all six instruments
and all four channels.
